\chapter{Continuous Mechanics of Rigid Bodies on Lie Groups}
\label{chapter:continous_mechanics_of_rigid_bodies_on_lie_groups}


This chapter introduces the continuous mechanics of rigid bodies on Lie groups, which is the foundation for the discrete mechanics formulation in Chapter~\ref{chapter:discrete_mechanics_for_rigid_bodies_on_lie_groups}. 
Therefore, Section \ref{sec:continuous_lagrangian_mechanics} presents the representation of geometric mechanics on Lie groups
by deriving the Euler-Lagrange equations and the Legendre transformation. Moreover, the properties of the Lagrangian flow are discussed. Consequently, 
Section \ref{section:continuous_3D_pendulum} applies the derived framework to the unforced 3D-pendulum as an example of a mechanical system on a Lie group, and verifies the momentum preservation properties of the Lagrangian 
flow of the model, based on the work of \cite{Lee2008}. 

% Chapter 3.1: Lagrangian Mechanics
% Required cross-reference in Chapter 2:
% Add \label{prop:infinitesimal_variation_matrix_lie_group}
% immediately after
% \begin{proposition}[Infinitesimal variation on a matrix Lie group]

\section{Lagrangian Mechanics}
\label{sec:continuous_lagrangian_mechanics}

Differential geometry provides a framework for describing geometric mechanics, by
exploring the geometric structure of a Lagrangian or Hamiltonian system \cite{Lee2008}.
Thereby, Lagrangian mechanics describes the motion of a mechanical system through
a variational principle. Instead of directly discretizing or integrating
the equations of motion, the trajectory is obtained from the stationary
value of the action integral. This formulation provides the continuous 
foundation for the discrete variational integrators introduced in 
Chapter~\ref{chapter:discrete_mechanics_for_rigid_bodies_on_lie_groups}. 
In particular, the variational construction keeps the symplectic and 
momentum-preserving structure of mechanical systems, which is later used 
for the geometric trajectory optimization formulation via Semidefinite Programming (SDP) Relaxation in Chapter~\ref{chapter:certifiable_trajectory_optimization}
\cite{Lee2008,Marsden2001,Sangli2025}.

\begin{figure}[ht]
    \centering
    \begin{tikzpicture}[
        node distance=0.9cm and 1.2cm,
        box/.style={
            draw, rectangle, align=center, minimum width=4.6cm,
            minimum height=1.2cm, inner sep=6pt
        },
        arrow/.style={-Latex}
    ]
        \node[box] (manifold){Configuration and velocity\\ $(\bm{g},\dot{\bm{g}})\in TG$};
        \node[box, below=of manifold] (lagrangian) {Lagrangian\\ $L(\bm{g},\bm{\xi})$};
        \node[box, below left=of lagrangian, xshift=1.3cm, yshift=-0.3cm] (action) {Action integral\\ $\mathcal S=\int_{t_0}^{t_f} L(\bm{g},\bm{\xi})\,dt$};
        \node[box, below right=of lagrangian, xshift=-1.3cm, yshift=-0.3cm] (legendre) {Legendre transform.\\ $\bm{\mu}=\mathbb{F}L(\bm{g},\bm{\xi})$};
        \node[box, below=of action] (variation) {Variation\\ $\delta\mathcal S=0$};
        \node[box, below=of legendre] (hamilton) {Hamilton equations};
        \node[box, below=of variation] (euler) {Euler--Lagrange equations};

        \draw[arrow] (manifold) -- (lagrangian);
        \draw[arrow] (lagrangian) -- (action);
        \draw[arrow] (lagrangian) -- (legendre);
        \draw[arrow] (action) -- (variation);
        \draw[arrow] (legendre) -- (hamilton);
        \draw[arrow] (variation) -- (euler);
    \end{tikzpicture}
    \caption{General method to derive the Euler-Lagrange equations of a mechanical system on a Lie group (adapted from \cite{Lee2008}).}
    \label{fig:continuous_mechanics}
\end{figure}

Figure~\ref{fig:continuous_mechanics}, adapted from \cite{Lee2008}, summarizes the general method to derive the Euler-Lagrange
equations of a mechanical system on a Lie group. Thereby, the path for which the action of a Lagrangian is stationary
over time has to be derived. The variation of the trajectory gives the Euler-Lagrange equations. 
Alternatively, the Legendre transformation can be used to derive the Hamiltonian equations of motion.   


\subsection{Euler--Lagrange Equations}
\label{sec:continuous_euler_lagrange_equations}

Following \cite{Lee2008}, the configuration manifold is a matrix Lie group $G$ with Lie algebra $\mathfrak g$. By left trivialization, the tangent bundle $TG$ is identified with $G\times\mathfrak g$ in \eqref{eq:left_trivialized_velocity}.
A mechanical system on $G$ is described by the Lagrangian
\begin{equation}
    L:G\times\mathfrak g\rightarrow\mathbb R,
    \qquad
    L(\bm{g},\bm{\xi})=T(\bm{g},\bm{\xi})-U(\bm{g}),
    \label{eq:continuous_lagrangian_lie_group}
\end{equation}
where $T$ and $U$ denote the kinetic and potential energy. On the interval $[t_0,t_f]$, the action integral is given by
\begin{equation}
    \mathcal S
    =
    \int_{t_0}^{t_f}L(\bm{g},\bm{\xi})\,\mathrm dt.
    \label{eq:continuous_action_integral}
\end{equation}
To derive the Euler--Lagrange equations, the variation of the action integral has to be computed. By definition, Hamilton's principle states that the action is stationary
\begin{equation}
    \delta\mathcal S
    =
    \delta\int_{t_0}^{t_f}L(\bm{g},\bm{\xi})\,\mathrm dt
    =
    \int_{t_0}^{t_f}\delta L(\bm{g},\bm{\xi})\,\mathrm dt
    =
    0,
    \label{eq:hamilton_principle_lie_group}
\end{equation}
for all variations with fixed endpoints at $t_0$ and $t_f$ \cite{Marsden2001}.
The variation of the Lagrangian is expressed as
\begin{equation}
    \delta L(\bm{g},\bm{\xi})
    =
    \left\langle D_{\bm{g}}L(\bm{g},\bm{\xi}),\delta\bm{g}\right\rangle
    +
    \left\langle D_{\bm{\xi}}L(\bm{g},\bm{\xi}),\delta\bm{\xi}\right\rangle,
    \label{eq:variation_continuous_lagrangian}
\end{equation}
with the partial derivatives $D_{\bm{g}}L(\bm{g},\bm{\xi})\in T_{\bm{g}}^*G$ with respect to $\bm{g}$ and $D_{\bm{\xi}}L(\bm{g},\bm{\xi})\in\mathfrak g^*$ with respect to $\bm{\xi}$ \cite{Marsden2001}.
\cite{Lee2008} extends the local group variation introduced in \eqref{eq:infinitesimal_variation_matrix_lie_group} to a trajectory defined by
\begin{equation}
    \bm{g}^{\epsilon}(t)
    =
    \bm{g}(t)\exp(\epsilon\bm{\eta}(t)),
    \qquad
    \text{for}
    \qquad
    \bm{\eta}(t_0)=\bm{\eta}(t_f)=\bm{0},
    \label{eq:continuous_curve_variation_lie_group}
\end{equation}
for a differentiable matrix-valued curve $\bm{\eta}:[t_0,t_f]\rightarrow\mathfrak g$ and defines the corresponding infinitesimal variations as
\begin{equation}
    \delta\bm{g}=\bm{g}\bm{\eta},
    \qquad
    \delta\bm{\xi}
    =
    \dot{\bm{\eta}}+\operatorname{ad}_{\bm{\xi}}\bm{\eta}.
    \label{eq:continuous_lie_group_variations}
\end{equation}
The endpoint conditions in \eqref{eq:continuous_curve_variation_lie_group} keep the initial and terminal configurations fixed and remove the boundary contribution arising from integration by parts.

Substituting \eqref{eq:variation_continuous_lagrangian} and \eqref{eq:continuous_lie_group_variations} into \eqref{eq:hamilton_principle_lie_group}, the variation of the action integral is given by
\begin{align}
    \delta\mathcal S
    =
    \int_{t_0}^{t_f}
    \Bigg\langle
        T_e^*\mathbb{L}_{\bm{g}}\,D_{\bm{g}}L(\bm{g},\bm{\xi})
        +\operatorname{ad}_{\bm{\xi}}^*D_{\bm{\xi}}L(\bm{g},\bm{\xi})
        -\frac{\mathrm d}{\mathrm dt}D_{\bm{\xi}}L(\bm{g},\bm{\xi}),
        \bm{\eta}
    \Bigg\rangle
    \mathrm dt
    =0,
    \label{eq:variation_action_after_integration_by_parts}
\end{align}
where $T_e^*\mathbb{L}_{\bm{g}}:T_{\bm{g}}^*G\rightarrow\mathfrak g^*$ denotes the dual of the differential $T_e\mathbb{L}_{\bm{g}}$ \cite{Lee2008}. Since \eqref{eq:variation_action_after_integration_by_parts} holds for all $\bm{\eta} \in \mathfrak{g}$ satisfying $\bm{\eta}(t_0)=\bm{\eta}(t_f)=\bm{0}$, the continuous Euler--Lagrange equations on $G$ are defined as
\begin{subequations}
\label{eq:continuous_euler_lagrange_lie_group}
\begin{align}
    \frac{\mathrm d}{\mathrm dt}D_{\bm{\xi}}L(\bm{g},\bm{\xi})
    -\operatorname{ad}_{\bm{\xi}}^*D_{\bm{\xi}}L(\bm{g},\bm{\xi})
    -T_e^*\mathbb{L}_{\bm{g}}\,D_{\bm{g}}L(\bm{g},\bm{\xi})
    &=0,
    \label{eq:continuous_euler_lagrange_momentum_balance}\\
    \dot{\bm{g}}&=\bm{g}\bm{\xi}.
    \label{eq:continuous_euler_lagrange_reconstruction}
\end{align}
\end{subequations}
The equation \eqref{eq:continuous_euler_lagrange_momentum_balance} is the dynamical representation of the system with the velocity $\bm{\xi}$ expressed in the body-fixed frame, while \eqref{eq:continuous_euler_lagrange_reconstruction} reconstructs the trajectory on the configuration manifold \cite{Lee2008}.

In order to derive the Euler-Lagrange equations in the following Sections, a general method for calculating the partial derivatives of the Lagrangian with respect to $\bm{\xi}$ and $\bm{g}$ has to be defined.

\begin{proposition}[Directional derivative representation of $D_{\bm{\xi}} L(\bm{g},\bm{\xi})$]
Let $G$ be a Lie group with Lie algebra $\mathfrak g$, and let $L:G\times\mathfrak g\to\mathbb R$
be smooth \cite{Hall2000LieGroups}. For any $\delta\bm{\xi}\in\mathfrak g$ the partial derivative $D_{\bm{\xi}} L(\bm{g},\bm{\xi})\in\mathfrak g^*$ is characterized by
\begin{equation}\label{eq:derivative_omega}
\langle D_{\bm{\xi}} L(\bm{g},\bm{\xi}),\delta\bm{\xi}\rangle
=\left.\frac{d}{d\varepsilon}\right|_{\varepsilon=0}L(\bm{g},\bm{\xi}+\varepsilon\,\delta\bm{\xi}).
\end{equation}
\end{proposition}

\begin{proof}
Since $\mathfrak g$ is a vector space, the differential of the map $L(\bm{g},\bm{\xi})$ at $\bm{\xi}$ is a linear map $\mathfrak g\to\mathbb R$, as $\bm{g}$ is fixed. Therefore the partial derivative $D_{\bm{\xi}}L(\bm{g},\bm{\xi})$ is an element of the dual space $\mathfrak g^*$. Moreover, for any $\delta\bm{\xi}\in\mathfrak g$, the curve
$\bm{\gamma}(\varepsilon)=\bm{\xi}+\varepsilon\,\delta\bm{\xi}$ satisfies $\bm{\gamma}(0)=\bm{\xi}$ and $\dot{\bm{\gamma}}(0)=\delta\bm{\xi}$, which realizes $\delta\bm{\xi}$ as a tangent direction at $\bm{\xi}$ for any $\varepsilon \in (-c,c)$ \cite{Lee2008}.
The corresponding directional derivative is
\[
\langle D_{\bm{\xi}} L(\bm{g},\bm{\xi}),\delta\bm{\xi}\rangle
=
\left.\frac{d}{d\varepsilon}\right|_{\varepsilon=0} L(\bm{g},\bm{\gamma}(\varepsilon))
=
\left.\frac{d}{d\varepsilon}\right|_{\varepsilon=0}L(\bm{g},\bm{\xi}+\varepsilon\,\delta\bm{\xi}),
\]
using $\left\langle D_{\bm{\xi}} L(\bm{g},\bm{\xi}),\delta\bm{\xi}\right\rangle
=
D_{\bm{\xi}} L(\bm{g},\bm{\xi})(\delta\bm{\xi})$.
\end{proof}


\begin{proposition}[Directional derivative representation of $D_{\bm{g}} L(\bm{g},\bm{\xi})$]
Let $G$ be a Lie group with Lie algebra $\mathfrak g$, and let $L:G\times\mathfrak g\to\mathbb R$
be smooth and $\bm{\eta}$ an arbitrary matrix-valued curve in $\mathfrak g$ \cite{Hall2000LieGroups}. For any $\delta\bm{g}=T_e\mathbb{L}_{\bm{g}}(\bm{\eta})=\bm{g}\bm{\eta}\in T_{\bm{g}}G$ \eqref{eq:continuous_lie_group_variations}
the partial derivative $D_{\bm{g}} L(\bm{g},\bm{\xi})\in T_{\bm{g}}^*G$ is characterized by
\begin{equation} \label{eq:derivative_rotation}
\langle D_{\bm{g}} L(\bm{g},\bm{\xi}),\delta\bm{g}\rangle
=
\langle T_{e}^{*}\mathbb{L}_{\bm{g}} D_{\bm{g}} L(\bm{g},\bm{\xi}),\bm{\eta} \rangle
=
\left.\frac{d}{d\varepsilon}\right|_{\varepsilon=0}L(\bm{g}\exp{(\varepsilon \bm{\eta})},\bm{\xi}).
\end{equation}
\end{proposition}

\begin{proof}
Since $G$ is a smooth manifold, the differential of the map $L(\bm{g},\bm{\xi})$ is a linear map $T_{\bm{g}}G\to\mathbb R$, as $\bm{\xi}$ is fixed. Therefore the partial derivative $D_{\bm{g}}L(\bm{g},\bm{\xi})$ is an element of the dual space $T_{\bm{g}}^*G$. 
Using the definition of the variation $\bm{g}^{\epsilon}$ \eqref{eq:continuous_curve_variation_lie_group}, the corresponding directional derivative is
\[
\langle D_{\bm{g}} L(\bm{g},\bm{\xi}),\delta\bm{g}\rangle
= \left.\frac{d}{d\varepsilon}\right|_{\varepsilon=0}L(\bm{g}^{\epsilon},\bm{\xi})
= \left.\frac{d}{d\varepsilon}\right|_{\varepsilon=0}L(\bm{g}\exp{(\varepsilon \bm{\eta})},\bm{\xi}),
\]
using $\left\langle D_{\bm{g}} L(\bm{g},\bm{\xi}),\delta\bm{g}\right\rangle
=
D_{\bm{g}} L(\bm{g},\bm{\xi})(\delta\bm{g})$ \cite{Lee2008}. Additionally, the directional derivative can be expressed by the covector $T_e^*\mathbb{L}_{\bm{g}}$ and
the definition of the
differential of the left translation $\bm{g}\bm{\eta} = T_e\mathbb{L}_{\bm{g}}(\bm{\eta})$ defined in \eqref{eq:preliminaries_differential_left_translation}, which leads to
\[
\langle D_{\bm{g}} L(\bm{g},\bm{\xi}),\delta\bm{g}\rangle
= \langle D_{\bm{g}} L(\bm{g},\bm{\xi}),\bm{g}\bm{\eta} \rangle
= \langle D_{\bm{g}} L(\bm{g},\bm{\xi}), T_e\mathbb{L}_{\bm{g}}(\bm{\eta}) \rangle
= \langle T_e ^* \mathbb{L}_{\bm{g}} D_{\bm{g}} L(\bm{g},\bm{\xi}), \bm{\eta} \rangle.
\]
\end{proof}

\subsection{Legendre Transformation and Hamiltonian Mechanics}
\label{sec:continuous_legendre_hamiltonian_mechanics}

In contrast to the Lagrangian $L:G\times\mathfrak{g}\to\mathbb{R}$ formulation, the Hamiltonian formulation $H: G \times \mathfrak{g}^\star \to \mathbb{R}$ represents the state on the cotangent bundle $T^*G$ identified with $G \times \mathfrak{g}^*$ by left trivialization. 
Thereby, the state is written as $(\bm{g},\bm{\mu})$ with the body-fixed momentum $\bm{\mu}\in\mathfrak g^*$ \cite{Marsden2001}. 
The Lagrangian and Hamiltonian descriptions are related by the Legendre transformation $\mathbb{F}L:G\times\mathfrak g\rightarrow G\times\mathfrak g^*$
\begin{equation}
    \mathbb{F}L:
    (\bm{g},\bm{\xi})
    \mapsto
    (\bm{g},\bm{\mu})
    =
    \left(
        \bm{g},
        D_{\bm{\xi}}L(\bm{g},\bm{\xi})
    \right),
    \label{eq:continuous_generalized_momentum}
\end{equation}
Consequently, applying the Legendre transformation~\eqref{eq:continuous_generalized_momentum} to the Euler-Lagrange equations \eqref{eq:continuous_euler_lagrange_lie_group}, the Hamiltonian equations of motion are derived as
\begin{subequations}
\label{eq:continuous_hamilton_equations_lie_group}
\begin{align}
    \dot{\bm{\mu}}
    -\operatorname{ad}_{\bm{\xi}}^*\bm{\mu}
    -T_e^*\mathbb{L}_{\bm{g}}\,D_{\bm{g}}L(\bm{g},\bm{\xi})
    &=0,
    \label{eq:continuous_hamilton_momentum_equation}\\
    \dot{\bm{g}}&=\bm{g}\bm{\xi}.
    \label{eq:continuous_hamilton_reconstruction}
\end{align}
\end{subequations}

\subsection{Properties of the Lagrangian Flow}
\label{sec:properties_continuous_lagrangian_flow}

\cite{Marsden2001} defines the Lagrangian flow $\mathcal{F}_L^T:G\times\mathfrak g\rightarrow G\times\mathfrak g$,
which maps an initial state to the corresponding state at time $t$. 
Thereby, for a regular Lagrangian, the Lagrangian flow evolves the system in time while preserving its geometric structure.
Under a group action, it has symplecticity and momentum-preservation properties \cite{Lee2008,Marsden2001}.

\paragraph{Symplecticity}
The Lagrangian one-form $\Theta_L$ on $G\times\mathfrak g$ is defined by
\begin{equation}
    \Theta_L(\bm{g},\bm{\xi})\cdot(\delta\bm{g},\delta\bm{\xi})
    =
    \left\langle
        D_{\bm{\xi}}L(\bm{g},\bm{\xi}),
        \bm{g}^{-1}\delta\bm{g}
    \right\rangle.
    \label{eq:continuous_lagrangian_one_form}
\end{equation}
With the Lagrangian symplectic two-form $\Omega_L=\mathrm d\Theta_L,$
the flow preserves symplecticity according to
\begin{equation}
    \left(F_L^T\right)^*\Omega_L=\Omega_L,
    \label{eq:continuous_lagrangian_flow_symplecticity}
\end{equation}
for $T=t_f-t_0$ \cite{Lee2008}.

\paragraph{Noether's theorem}
Let a left action of a Lie group $H$ with Lie algebra $\mathfrak h$ act  
smoothly on $G$ by the mapping $\Phi:H\times G\rightarrow G$ such that
$\Phi_{\bm{h}}(\bm{g})=\Phi(\bm{h},\bm{g})$. \cite{Lee2008} defines the left trivialization 
$\phi_L:TG \to G\times\mathfrak g$ by
\begin{equation}\label{eq:left_trivialization_noethers_theorem}
    \phi_L(\bm{g},\dot{\bm{g}}) = (\bm{g}, \bm{g}^{-1}\dot{\bm{g}}),
    \qquad
    \phi_L^{-1}(\bm{g},\bm{g}^{-1}\dot{\bm{g}}) = (\bm{g},\dot{\bm{g}}).
\end{equation}
Additionally, the infinitesimal generator $\zeta_{G\times\mathfrak{g}}: G \times \mathfrak g \to T(G\times\mathfrak g)$ is given by
\begin{equation}
    \zeta_{G\times\mathfrak g}(\bm{g},\bm{\xi})
    =
    \left.\frac{\mathrm d}{\mathrm d\epsilon}\right|_{\epsilon=0}
        \phi_L
        \circ T_{\bm{g}}\Phi_{\exp_H(\epsilon\zeta)}(\bm{g})
        \cdot (\phi_L^{-1}(\bm{g},\bm{\xi})),
    \label{eq:lifted_infinitesimal_generator}
\end{equation}
for $\zeta\in\mathfrak h$ \cite{Lee2008}.
Moreover, the corresponding Lagrangian momentum map $J_L:G\times\mathfrak g\rightarrow\mathfrak h^*$ is defined by
\begin{equation}
    J_L(\bm{g},\bm{\xi})\cdot\zeta
    =
    \Theta_L\cdot
    \zeta_{G\times\mathfrak g}.
    \label{eq:continuous_lagrangian_momentum_map}
\end{equation}

\begin{theorem}[Noether's theorem]
\label{thm:continuous_noether}
If the Lagrangian is infinitesimally invariant under the lifted action, such that
\begin{equation}
    \mathrm dL(\bm{g},\bm{\xi})\cdot
    \zeta_{G\times\mathfrak g}(\bm{g},\bm{\xi})=0
    \qquad
    \text{for all }\zeta\in\mathfrak h,
    \label{eq:infinitesimal_lagrangian_invariance}
\end{equation}
then the associated momentum map is preserved along the Lagrangian flow,
\begin{equation}
    J_L\!\left(F_L^t(\bm{g},\bm{\xi})\right)
    =
    J_L(\bm{g},\bm{\xi}),
    \label{eq:continuous_noether_momentum_preservation}
\end{equation}
which is considered as Noether's theorem \cite{Marsden2001}.
\end{theorem}

Noether's theorem connects a continuous symmetry of the Lagrangian with a conserved momentum quantity \cite{Lee2008,Marsden2001}. 
This result is applied in Section~\ref{sec:continuous_3D_pendulum_momentum_validation} to identify and verify the 
momentum preserved by the unforced 3D-pendulum. However, \cite{Marsden2001} states that 
the variational integrator does not generally conserve the continuous energy of the system exactly.
Nevertheless, due to the symplectic structure, the energy is bounded and oscillates around the exact energy over a long-time horizon.


\section{Unforced 3D-Pendulum}
\label{section:continuous_3D_pendulum}

In the following Section, the unforced mechanics of a three-degree-of-freedom pendulum given by \cite{Lee2008} are presented, using Lagrangian mechanics of Section \ref{sec:continuous_lagrangian_mechanics}. Following, the momentum preservation properties of the Lagrangian flow of the model will be verified. 

\subsection{Derivation of Lie Group Mechanics}

\paragraph{Configuration Manifold}

\begin{figure}[t!]
	\centering
	\includegraphics[width=0.75\textwidth]{pendulum_3d_so3.pdf}
	\caption{3D-pendulum system attached to a fixed pivot.}
	\label{fig:3d_pendulum}
\end{figure}

The 3D-Pendulum model, illustrated in \ref{fig:3d_pendulum}, consist of a rigid body rotating under the influence of gravity $g$ attached to a fixed frictionless pivot \cite{Shen20043DPendulum}. Moreover, the system has no external control moment, since it is unforced. A space frame and a body-fixed frame attached to the pendulum, both with their origins at the pivot, are defined. The unit vector $\bm{e}_3$ represents the gravity direction in the space frame, while $\bm{\rho_c}$ denotes the vector from the pivot to the center of mass expressed in the body-fixed frame. Following the fixed pivot, there is no translational degree of freedom. Therefore, the configuration manifold of the model is given by the special orthogonal group
\[
SO(3)
=
\left\{
\bm{R}\in\mathbb{R}^{3\times 3}
\;\middle|\;
\bm{R}^\top \bm{R}=\bm{I}_3,\ 
\det(\bm{R})=1
\right\}.\]
The configuration is completely determined by the rotation matrix $\bm{R}\in SO(3)$, which maps the body-fixed frame to the space frame. Additionally, in $SO(3)$ the attitude kinematics equation is defined by 
\begin{equation}
\label{eq:SO3_attitude_kinematics}
\dot{\bm{R}} = \bm{R}\hat{\bm{\omega}},
\end{equation}
with the angular velocity $\bm{\omega}\in \mathbb{R}^3$, mapped to the Lie algebra $\mathfrak{so}(3)$ by the hat map $\hat{\cdot}: \mathbb{R}^3 \to \mathfrak{so}(3)$ \cite{Lee2008}. Consequently, $\hat{\bm{\omega}}$ is a $3 \times 3$ skew-symmetric matrix. 

\paragraph{Lagrangian}

To derive the Lagrangian $L:SO(3) \times\mathfrak{so}(3) \to \mathbb{R}$, the kinetic $T:SO(3) \times\mathfrak{so}(3) \to \mathbb{R} $ and potential $V:SO(3) \to \mathbb{R}$ energy of the rigid body are introduced. The rotational inertia of the mass is expressed in the body-frame by the inertia matrix $\bm{J} \in \mathbb{R}^{3 \times 3}$. \cite{Lee2008} defines 
\begin{equation}\label{eq:3D_pendulum_kinetic_energy}
T(\bm{\omega}) = \frac{1}{2} \int_{B} \lVert \hat{\bm{\omega} \bm{\rho}} \rVert^{2} \, dm(\bm{\rho})
= \frac{1}{2}\operatorname{tr}[\hat{\bm{\omega}} \bm{J}_d \hat{\bm{\omega}}^\top]
= \frac{1}{2} \bm{\omega}^\top \bm{J} \bm{\omega},
\end{equation}
as a general expression for the kinetic energy of the single body $B$ with the velocity $\bm{\omega} \times \bm{\rho} = \hat{\bm{\omega}}\bm{\rho}$. Thereby, 
\begin{equation}\label{eq:non-standard_inertia}
    \bm{J}_d = \int_{B} \hat{\bm{\rho}} \hat{\bm{\rho}}^\top  \, dm(\bm{\rho})
    =\frac{1}{2}\operatorname{tr}[\bm{J}]\,\bm{I}-\bm{J},
\end{equation}
is a nonstandard inertia matrix, with the relation
\begin{equation} \label{eq:nonstandard_Inertia_Relation}
    (\bm{J}\bm{\omega})^\wedge = \hat{\bm{\omega}}\bm{J}_d + \bm{J}_d \hat{\bm{\omega}},
\end{equation}
for any $\bm{\omega}\in\mathbb{R}^3$. The potential energy is given by
\begin{equation}\label{eq:3D_pendulum_potential_energy}
    V(\bm{R}) = -mg\, \bm{e}_3^\top \bm{R}\bm{\rho}_c,
\end{equation}
where \(\bm{R}\bm{\rho}_c\) denotes the vector from the pivot to the center of mass expressed in the space-fixed frame, and \(\bm{e}_3^\top \bm{R}\bm{\rho}_c\) is its projection onto the gravity direction \(\bm{e}_3\) \cite{Lee2008}. Using \eqref{eq:3D_pendulum_kinetic_energy} and \eqref{eq:3D_pendulum_potential_energy}, the Lagrangian of the unforced 3D pendulum is given by 
\begin{equation}\label{eq:continuous_3D_pendulum_lagrangian}
    L(\bm{R}, \hat{\bm{\omega}}) = T(\hat{\bm{\omega}}) - V(\bm{R})
    = \frac{1}{2}\operatorname{tr}[\hat{\bm{\omega}} \bm{J}_d \hat{\bm{\omega}}^\top] + mg\, \bm{e}_3^\top \bm{R}\bm{\rho}_c.
\end{equation}


\paragraph{Euler--Lagrange Equations}

In order to define the Euler--Lagrange equations \eqref{eq:continuous_euler_lagrange_lie_group}, the derivatives of \eqref{eq:continuous_3D_pendulum_lagrangian} have to be calculated. Firstly, the directional derivative $D_{\hat{\bm{\omega}}}L(\bm{R},\hat{\bm{\omega}})$ is given by applying \eqref{eq:derivative_omega}, which yields
\begin{align}
\left\langle D_{\hat{\bm{\omega}}}L(\bm{R},\hat{\bm{\omega}}),\,\delta\hat{\bm{\omega}}\right\rangle
&=
\left.\frac{d}{d\epsilon}\right|_{\epsilon=0}
L\!\left(\bm{R},\hat{\bm{\omega}}+\epsilon\,\delta\hat{\bm{\omega}}\right) \notag\\
&=
\left.\frac{d}{d\epsilon}\right|_{\epsilon=0}
\frac{1}{2}\operatorname{tr}\!\left[
(\hat{\bm{\omega}}+\epsilon\,\delta\hat{\bm{\omega}})
\,\bm{J}_d\,
(\hat{\bm{\omega}}+\epsilon\,\delta\hat{\bm{\omega}})^{T}
\right] \notag\\
&=
\frac{1}{2}\operatorname{tr}\!\left[
\delta\hat{\bm{\omega}}\,\bm{J}_d\,\hat{\bm{\omega}}^{T}
+
\hat{\bm{\omega}}\,\bm{J}_d\,(\delta\hat{\bm{\omega}})^{T}
\right] \notag.
\end{align}
By using the antisymmetry identity of a cross product $\hat{\bm{x}}=-\hat{\bm{x}}^\top$ for any $\bm{x} \in \mathbb{R}^3$ and $\operatorname{tr}[\bm{A}\bm{B}]=\operatorname{tr}[\bm{B}\bm{A}]$ for any $\bm{A},\bm{B} \in \mathbb{R}^{3 \times 3}$, 
the directional derivative is rewritten as
\begin{align}\label{eq:continuous_3D_pendulum_derivative_omega}
\left\langle D_{\hat{\bm{\omega}}}L(\bm{R},\hat{\bm{\omega}}),\,\delta\hat{\bm{\omega}}\right\rangle
&=
\frac{1}{2}\operatorname{tr}\!\left[
-\delta\hat{\bm{\omega}}\,\bm{J}_d\,\hat{\bm{\omega}}
-
\hat{\bm{\omega}}\,\bm{J}_d\,\delta\hat{\bm{\omega}}
\right] \notag \\
&=
\frac{1}{2}\operatorname{tr}\!\left[
-\bigl(\bm{J}_d\hat{\bm{\omega}}+\hat{\bm{\omega}}\bm{J}_d\bigr)
\,\delta\hat{\bm{\omega}}
\right] \notag\\
&=
\left\langle \bm{J}_d\hat{\bm{\omega}}+\hat{\bm{\omega}}\bm{J}_d,\,
\delta\hat{\bm{\omega}}\right\rangle \notag\\
&=
\left\langle (\bm{J}\bm{\omega})^\wedge,\,\delta\hat{\bm{\omega}}\right\rangle,
\end{align}
with the definition of the inner product \eqref{eq:trace_inner_product} and the skew-symmetric property of \eqref{eq:nonstandard_Inertia_Relation}. 

\indent Moreover, the derivative $D_{\bm{R}}L(\bm{R},\hat{\bm{\omega}})$, which in this model depends only on the potential energy $V(\bm{R})$, is calculated using \eqref{eq:derivative_rotation} as
\begin{align}
\left\langle D_{\bm{R}} L(\bm{R},\hat{\bm{\omega}}),\,\delta \bm{R}\right\rangle
&=
\left\langle T_{e}^{*}\mathbb{L}_{\bm{R}}\,D_{\bm{R}} L(\bm{R},\hat{\bm{\omega}}),\, \hat{\bm{\eta}} \right\rangle \notag\\
&=
\left.\frac{d}{d\epsilon}\right|_{\epsilon=0}
L\!\left(\bm{R}\exp(\epsilon\hat{\bm{\eta}}),\hat{\bm{\omega}}\right) \notag\\
&=
\left.\frac{d}{d\epsilon}\right|_{\epsilon=0}
mg\,\bm{e}_3^{\top}\bm{R}\exp(\epsilon\hat{\bm{\eta}})\bm{\rho}_c \notag\\
&=
mg\,\bm{e}_3^{\top}\bm{R}\hat{\bm{\eta}}\,\bm{\rho}_c \notag\\
&=
mg\,(\bm{R}^{\top}\bm{e}_3)^{\top}\hat{\bm{\eta}}\,\bm{\rho}_c \notag\\
&=
mg\,(\bm{\rho}_c\times \bm{R}^{\top}\bm{e}_3)\cdot \bm{\eta}, \notag
\end{align}
since $\bm{a}^\top \hat{\bm{x}}\bm{b} = \bm{a} \cdot(\bm{x} \times \bm{b}) = (\bm{b} \times \bm{a}) \cdot \bm{x}$ for any $\bm{a}, \bm{b}, \bm{x} \in \mathbb{R}^3$. Applying the definition of inner products \eqref{eq:trace_inner_product}, it follows that 
\begin{equation}\label{eq:continuous_3D_pendulum_derivative_rotation}
    \left\langle T_{e}^{*}\mathbb{L}_{\bm{R}}\,D_{\bm{R}} L(\bm{R},\hat{\bm{\omega}}),\, \hat{\bm{\eta}} \right\rangle
    = \left\langle (mg\,(\bm{\rho}_c\times \bm{R}^{\top}\bm{e}_3))^\wedge, \hat{\bm{\eta}}\right\rangle.
\end{equation}

Since the directional derivative
\[
D_{\hat{\bm{\omega}}}L(\bm{R},\hat{\bm{\omega}})
\;\widehat{=}\;
\bm{J}_d\hat{\bm{\omega}}+\hat{\bm{\omega}}\bm{J}_d,
\]
was derived in \eqref{eq:continuous_3D_pendulum_derivative_omega}, the coadjoint term in the Euler--Lagrange equations can be written in matrix form by using the identity \eqref{eq:ad_lie_bracket}. Therefore,
\begin{align}
ad^*_{\hat{\bm{\omega}}}D_{\hat{\bm{\omega}}}L(\bm{R},\hat{\bm{\omega}})
&=
\left[\bm{J}_d\hat{\bm{\omega}}+\hat{\bm{\omega}}\bm{J}_d,\hat{\bm{\omega}}\right] \notag \\
&=
(\bm{J}_d\hat{\bm{\omega}}+\hat{\bm{\omega}}\bm{J}_d)\hat{\bm{\omega}}
-
\hat{\bm{\omega}}(\bm{J}_d\hat{\bm{\omega}}+\hat{\bm{\omega}}\bm{J}_d) \notag \\
&=
(\bm{J}\bm{\omega} \times \bm{\omega})^\wedge \notag
\end{align}
with the relation $\hat{\bm{x}}\hat{\bm{y}}-\hat{\bm{y}}\hat{\bm{x}}=(\bm{x} \times \bm{y})^\wedge$ for any $\bm{x},\bm{y} \in \mathbb{R}^3$ and \eqref{eq:nonstandard_Inertia_Relation}. Furthermore, by the antisymmetry of the cross product, it follows that
\begin{equation} \label{eq:continuous_3D_pendulum_adjoint_derivative_omega}   
ad^*_{\hat{\bm{\omega}}}D_{\hat{\bm{\omega}}}L(\bm{R},\hat{\bm{\omega}})
= (\bm{J}\bm{\omega} \times \bm{\omega})^\wedge
= -(\bm{\omega} \times \bm{J} \bm{\omega})^\wedge
\end{equation}
\indent Finally, the Euler-Lagrange equations of the unforced 3D-pendulum are expressed in $\mathfrak{so}(3)$ and $SO(3)$ form, by substituting \eqref{eq:continuous_3D_pendulum_derivative_omega}, \eqref{eq:continuous_3D_pendulum_derivative_rotation} and \eqref{eq:continuous_3D_pendulum_adjoint_derivative_omega} into \eqref{eq:continuous_euler_lagrange_lie_group}, which results in
\begin{subequations}\label{eq:continuous_3D_pendulum_euler_lagrange}
\begin{align}
 (\bm{J}\dot{\bm{\omega}})^\wedge + (\bm{\omega}\times (\bm{J}\bm{\omega}))^\wedge
 &= (mg\,\bm{\rho}_c \times \bm{R}^T\bm{e}_3)^\wedge,
 \label{eq:continuous_3D_pendulum_euler_lagrange_dynamics} \\
 \dot{\bm{R}} &= \bm{R}\hat{\bm{\omega}}.
 \label{eq:continuous_3D_pendulum_euler_lagrange_kinematics}
\end{align}
\end{subequations}

\paragraph{Legendre Transformation}

In order to obtain Hamilton's equations for the unforced 3D-pendulum model, the Legendre transformation $\mathbb{F}L:(SO(3)\times\mathfrak{so}(3)) \to (SO(3) \times \mathfrak{so}(3)^*)$ \eqref{eq:continuous_generalized_momentum} has to be applied, which results in the expression for the momentum
\begin{equation}  \label{eq:continuous_3D_pendulum_momentum}
    \hat{\bm{\Pi}} = D_{\hat{\bm{\omega}}}L(\bm{R},\hat{\bm{\omega}}) = (\bm{J}\bm{\omega})^\wedge
\end{equation}
where $ \hat{\bm{\Pi}} \in \mathfrak{so}(3)^*$. By substituting \eqref{eq:continuous_3D_pendulum_momentum} into \eqref{eq:continuous_3D_pendulum_euler_lagrange}, Hamilton's equations of an unforced 3D-pendulum are 
\begin{subequations}\label{eq:continuous_3D_pendulum_hamilton}
\begin{align}
    \hat{\dot{\bm{\Pi}}} + (\bm{J}^{-1}\bm{\Pi} \times \bm{\Pi})^\wedge &= (mg\,\bm{\rho}_c \times \bm{R}^T \bm{e}_3)^\wedge,
    \label{eq:continuous_3D_pendulum_hamilton_momentum} \\
    \dot{\bm{R}} &= \bm{R}(\bm{J}^{-1}\bm{\Pi})^\wedge,
    \label{eq:continuous_3D_pendulum_hamilton_kinematics}
\end{align}
\end{subequations}
with $\bm{J}$ being nonsingular.
Finally, the derived Euler-Lagrange equations \eqref{eq:continuous_3D_pendulum_euler_lagrange} and Hamilton's equations \eqref{eq:continuous_3D_pendulum_hamilton} correspond to the given 3D-pendulum results of \cite{Lee2008}.

\subsection{Verification of Momentum Conservation}
\label{sec:continuous_3D_pendulum_momentum_validation}


The verification of momentum conservation is particularly relevant,
in order to ensure that the numerical integration of the system's dynamics preserves the physical properties of the system over long time. 
This is essential for accurate numerical simulation and for 
reliable trajectory optimization methods \cite{Sangli2025}. Therefore, in the following, the momentum conservation property of the continuous 3D pendulum is verified based on its symmetry and Noether's theorem \ref{thm:continuous_noether}.

\cite{Lee2008} identifies a symmetry of the unforced 3D-pendulum's Lagrangian, which is infinitesimally invariant under an action of the group $H = SO(2) \simeq S^1$. Here, a rotation about the gravity direction $\bm{e}_3$ does not change the Lagrangian. Thus, $SO(2)$ and $S^1$ describe the same symmetry group under the left action $\Phi: S^1 \times SO(3) \to SO(3)$ of $H$ on $G$,
\begin{equation}\label{eq:continuous_3D_pendulum_symmetry_action}
    \Phi(\theta, \bm{R})= \exp_{SO(3)}(\theta \hat{\bm{e}_3})\bm{R},
\end{equation}
where $\exp_{SO(3)}(\theta \hat{\bm{e}_3})$ expresses a finite rotation about the gravity direction $\bm{e}_3$ by the angle $\theta$ \cite{Lee2008}. By Noether's Theorem \ref{thm:continuous_noether}, the infinitesimal invariance of the 3D-pendulum's Lagrangian is verified, if 
\begin{equation}\label{eq:continuous_3D_momentum_conservation}
    dL(\bm{R},\hat{\bm{\omega}})\cdot \zeta_{SO(3)\times\mathfrak{so}(3)}(\bm{R}, \hat{\bm{\omega}}) = 0.
\end{equation}
Consequently, preserving the Lagrangian momentum map along the flow $J_L\!\left(\mathcal{F}_L^T(\bm{R},\hat{\bm{\omega}})\right)=J_L(\bm{R},\hat{\bm{\omega}})$.

Therefore, the infinitesimal generator of the lifted group action on $SO(3)\times\mathfrak{so}(3)$ is required \cite{Lee2008}. With the attitude kinematics relation \eqref{eq:SO3_attitude_kinematics}, the left trivialization $\phi_L(\bm{R},\dot{\bm{R}})=\phi_L(\bm{R},\bm{R}\hat{\bm{\omega}})=(\bm{R},T_e\mathbb{L}_{\bm{R}^{-1}}\bm{R}\hat{\bm{\omega}})=(\bm{R},\hat{\bm{\omega}})$, and its inverse $\phi_L^{-1}(\bm{R},\hat{\bm{\omega}})=(\bm{R},T_e\mathbb{L}_{\bm{R}}\hat{\bm{\omega}})=(\bm{R},\bm{R}\hat{\bm{\omega}})$ are taken from \eqref{eq:left_trivialization_noethers_theorem}, with the differential $T_{\bm{R}}$ of \eqref{eq:preliminaries_differential_left_SO}. Following \cite{Lee2008}, the lifted infinitesimal generator $\zeta_{SO(3) \times \mathfrak{so}(3)}:SO(3) \times \mathfrak{so}(3) \to T(SO(3) \times \mathfrak{so}(3))$ \eqref{eq:lifted_infinitesimal_generator} is computed as
\begin{align}
\zeta_{SO(3)\times\mathfrak{so}(3)}(\bm{R},\hat{\bm{\omega}})
&=
\left.\frac{d}{d\epsilon}\right|_{\epsilon=0}
\phi_L \circ T_{\bm{R}}\Phi_{\exp_{SO(2)}(\epsilon\zeta)}
\circ \phi_L^{-1}(\bm{R},\hat{\bm{\omega}}) \notag\\
&=
\left.\frac{d}{d\epsilon}\right|_{\epsilon=0}
\phi_L \circ T_{\bm{R}}\Phi_{\exp_{SO(2)}(\epsilon\zeta)}
(\bm{R},\bm{R}\hat{\bm{\omega}}) \notag\\
&=
\left.\frac{d}{d\epsilon}\right|_{\epsilon=0}
\phi_L\!\left(
\exp(\epsilon\zeta\hat{\bm{e}}_3)\bm{R},
\exp(\epsilon\zeta\hat{\bm{e}}_3)\bm{R}\hat{\bm{\omega}}
\right) \notag\\
&=
\left.\frac{d}{d\epsilon}\right|_{\epsilon=0}
\left(
\exp(\epsilon\zeta\hat{\bm{e}}_3)\bm{R},
\bm{R}^{-1}\exp(\epsilon\zeta\hat{\bm{e}}_3)^{-1}
\bigl(\exp(\epsilon\zeta\hat{\bm{e}}_3)\bm{R}\hat{\bm{\omega}}\bigr)
\right) \notag\\
&=
\left.\frac{d}{d\epsilon}\right|_{\epsilon=0}
\left(
\exp(\epsilon\zeta\hat{\bm{e}}_3)\bm{R},
\hat{\bm{\omega}}
\right) \notag\\
&=
\left(
\zeta\hat{\bm{e}}_3\bm{R},
\bm{0}
\right).
\label{eq:continuous_3D_pendulum_lifted_generator}
\end{align}

This result can be substituted into \eqref{eq:continuous_3D_momentum_conservation}, using \eqref{eq:continuous_3D_pendulum_derivative_rotation}, which gives
\begin{align}
dL(\bm{R},\hat{\bm{\omega}})\cdot
\zeta_{SO(3)\times\mathfrak{so}(3)}(\bm{R},\hat{\bm{\omega}})
&=
dL(\bm{R},\hat{\bm{\omega}})\cdot
(\zeta \hat{\bm{e}}_3\bm{R},\bm{0}) \notag\\
&=
\left\langle D_{\bm{R}}L(\bm{R},\hat{\bm{\omega}}),\,\zeta \hat{\bm{e}}_3\bm{R}\right\rangle
+
\left\langle D_{\hat{\bm{\omega}}}L(\bm{R},\hat{\bm{\omega}}),\,\bm{0}\right\rangle \notag\\
&=
\left\langle T_e^*\mathbb{L}_{\bm{R}}D_{\bm{R}}L(\bm{R},\hat{\bm{\omega}}),\,
\hat{\bm{\eta}}\right\rangle.
\label{eq:continuous_3D_pendulum_infinitesimal_invariance_start}
\end{align}
Since the lifted generator has no component in the $\hat{\bm{\omega}}$ direction, the component directional derivative $D_{\hat{\bm{\omega}}}$ equals zero in the zero direction. Moreover, by the definition of the lifted infinitesimal generator \eqref{eq:lifted_infinitesimal_generator}, the result of \eqref{eq:continuous_3D_pendulum_lifted_generator} redefines the variational relationship as
\begin{equation} \label{eq:continuous_3D_pendulum_lifted_action_generator_variation}
    (\delta \bm{R}, \delta \hat{\bm{\omega}}) = (\zeta \hat{\bm{e}}_3\bm{R},\bm{0}).
\end{equation}
Consequently, using the previously given infinitesimal variation \eqref{eq:continuous_lie_group_variations}, the corresponding relation $\delta \bm{R} = \zeta \hat{\bm{e}}_3 \bm{R} = \bm{R} \hat{\bm{\eta}}$ is defined, from which the variation $\hat{\bm{\eta}}$ can be expressed as
\begin{equation}\label{eq:continuous_3D_pendulum_variation_eta}
    \hat{\bm{\eta}}= \zeta \bm{R}^{-1} \hat{\bm{e}}_3 \bm{R}.
\end{equation}
This can be substituted into \eqref{eq:continuous_3D_pendulum_infinitesimal_invariance_start}, yielding
\begin{align}
\left\langle T_e^*\mathbb{L}_{\bm{R}}D_{\bm{R}}L(\bm{R},\hat{\bm{\omega}}),\,
\zeta \bm{R}^{-1}\hat{\bm{e}}_3\bm{R}\right\rangle
&=
\left.\frac{d}{d\epsilon}\right|_{\epsilon=0}
L\!\left(\bm{R}\exp\!\left(\epsilon \zeta \bm{R}^{-1}\hat{\bm{e}}_3\bm{R}\right),\hat{\bm{\omega}}\right).
\notag
\end{align}
Since only the potential term of the 3D-pendulum $V(\bm{R})$ \eqref{eq:3D_pendulum_potential_energy} depends on $\bm{R}$ and \eqref{eq:continuous_3D_pendulum_variation_eta} is skew symmetric, it can be shown that 
\begin{align} \label{eq:continuous_3d_pendulum_validation}
\left.\frac{d}{d\epsilon}\right|_{\epsilon=0}
L\!\left(\bm{R}\exp\!\left(\epsilon \zeta \bm{R}^{-1}\hat{\bm{e}}_3\bm{R}\right),\hat{\bm{\omega}}\right)
&=
mg\,\bm{e}_3^T\bm{R}
\left(\zeta \bm{R}^{-1}\hat{\bm{e}}_3\bm{R}\right)
\bm{\rho}_c \notag\\
&=
mg\,(\bm{\rho}_c^\top
\left(\zeta \bm{R}^{-1}\hat{\bm{e}}_3\bm{R}\right)^\top
\bm{R}^\top\bm{e}_3)^\top \notag\\
&=
-\,mg\,\bm{\rho}_c^\top
\left(\zeta \bm{R}^{-1}\hat{\bm{e}}_3\bm{R}\right)
\bm{R}^\top\bm{e}_3 \notag\\
&=
-\,mg\,\bm{\rho}_c^\top
\zeta \bm{R}^{-1}\hat{\bm{e}}_3\bm{e}_3 \notag\\
&=
-\,mg\,\bm{\rho}_c^\top
\zeta \bm{R}^{-1}(\bm{e}_3\times \bm{e}_3) \notag\\
&=0,
\end{align}
because $\bm{e}_3\times \bm{e}_3=\bm{0}$ and $\bm{R}^{-1}=\bm{R}^{\top}$ for any $\bm{R} \in SO(3)$. Therefore, the condition of \eqref{eq:continuous_3D_momentum_conservation} is verified and the Lagrangian is infinitesimally invariant under the symmetry action. Consequently, by Noether's Theorem ~\ref{thm:continuous_noether}, 
the corresponding momentum map $J_L(\bm{R},\bm{\omega})$ of the continuous 3D-pendulum is preserved along the Lagrangian flow
\[  J_L\!\left(\mathcal{F}_L^T(\bm{R},\hat{\bm{\omega}})\right)
    =
    J_L(\bm{R},\hat{\bm{\omega}}).
\]
